
c'est un anneau d'intersection complète (EGA IV 19.3.2), et par suite $f$ est d'intersection complète en $x$ au sens de EGA IV 19.3.6.

Réciproquement, supposons que $f$ soit un morphisme d'intersection complète au sens de EGA IV 19.3.6, et soit $x\in X$. Soit $U$ un voisinage affine de $x$, au-dessus d'un voisinage affine $U'$ de $y=f(x)$. On peut alors trouver un $U'$-schéma lisse $V$, et une immersion fermée $U\xrightarrow{i}V$ (en prenant par exemple pour $V$ un fibré vectoriel sur $U'$). Nous allons montrer que $i$ est régulière. Or d'après EGA IV 19.3.7, $i$ est transversalement régulier en $x$, puisque $V$ est lisse sur $Y$, donc $V_y$ régulier en $x$, et $U'$ plat sur $Y$. Il résulte alors de EGA IV 19.2.4, a)$\Rightarrow$c), que $i$ est régulière en $x$. Par suite, $f$ est un morphisme d'intersection complète au sens de 1.1.

\textbf{Proposition 1.5.} \uline{Soient} $f:X\to Y$, \uline{et} $g:Y\to Z$ \uline{deux morphismes localement d'intersection complète}. \uline{Alors} $g\circ f$ \uline{est un morphisme localement d'intersection complète}.

Soit $x\in X$. On peut trouver des voisinages ouverts $U,V$ de $x,f(x)$, et des immersions régulières $i,j$ telles qu'on ait les factorisations :
\[
\begin{tikzcd}[column sep=large,row sep=large]
U \arrow[r,"i"] \arrow[dr,"f"'] & V[T_1,\ldots,T_n] \arrow[d] &\\
& V \arrow[r,"j"] \arrow[dr,"g"'] & Z[T'_1,\ldots,T'_m] \arrow[d]\\
& & Z .
\end{tikzcd}
\]
On peut compléter ce diagramme par le carré commutatif :
